\documentclass[10pt]{beamer}

% --- Configurazione Tema ---
\usetheme{Madrid}
\usecolortheme{default}

% --- Pacchetti ---
\usepackage[utf8]{inputenc}
\usepackage[T1]{fontenc}
\usepackage[italian]{babel}
\usepackage{tikz}
\usetikzlibrary{arrows.meta, positioning, shapes.geometric, decorations.pathreplacing}
\usepackage{listings}
\usepackage{xcolor}

% --- Colori Personalizzati ---
\definecolor{primary}{RGB}{38, 66, 139}
\setbeamercolor{structure}{fg=primary}
\setbeamercolor{title}{fg=white, bg=primary}

% --- Configurazione Codice ---
\lstset{
    basicstyle=\ttfamily\scriptsize,
    breaklines=true,
    frame=single,
    backgroundcolor=\color{gray!10},
    keywordstyle=\color{blue},
    commentstyle=\color{green!50!black},
    stringstyle=\color{orange}
}

% --- Informazioni ---
\title[Web Statico vs Dinamico]{Programmazione e Tecnologie Web}
\subtitle{Lezione 2: Web Statico, Dinamico e Logica Client-Side}
\author{Giuseppe Capasso}
\institute[ITS]{Istituto Tecnico Superiore (ITS)}
\date{A.A. 2025/2026}

\begin{document}

% --- Slide Titolo ---
\begin{frame}
    \titlepage
\end{frame}

% --- Obiettivi ---
\begin{frame}{Obiettivi della Lezione}
    \begin{itemize}
        \item \textbf{Distinguere tra Web Statico e Dinamico}
        \item \textbf{Paradigmi di Rendering}: SSR vs CSR
        \item \textbf{Static Site Generation}: Focus su Hugo
        \item \textbf{Deployment Moderno}: Cloudflare Pages e CDN
        \item \textbf{Persistenza}: localStorage e paradigma CRUD
        \item \textbf{Laboratorio}: Todo App con jQuery
    \end{itemize}
\end{frame}

% --- Web Statico ---
\section{Web Statico vs Dinamico}
\begin{frame}{Il Web Statico: Caratteristiche}
    \begin{columns}
        \begin{column}{0.5\textwidth}
            \textbf{Cos'è?}
            \begin{itemize}
                \item File pre-esistenti sul server.
                \item Nessun database coinvolto.
                \item Risposta istantanea.
            \end{itemize}
            \vspace{0.5cm}
            \textbf{Vantaggi:}
            \begin{itemize}
                \item Performance estreme.
                \item Massima sicurezza.
                \item Costi di hosting minimi.
            \end{itemize}
        \end{column}
        \begin{column}{0.5\textwidth}
            \centering
            \begin{tikzpicture}[scale=0.6, transform shape]
                \node (client) [draw, rounded corners, fill=primary!10, minimum width=2cm] {Client};
                \node (server) [draw, rounded corners, fill=primary!10, right=2cm of client, minimum width=2cm] {Server Web};
                \node (disk) [draw, shape=cylinder, aspect=0.5, rotate=90, fill=gray!20, below=1cm of server] {\rotatebox{-90}{Files}};
                \draw [->, >=stealth, primary] ([yshift=0.2cm]client.east) -- node[above] {GET} ([yshift=0.2cm]server.west);
                \draw [<-, >=stealth, primary] ([yshift=-0.2cm]client.east) -- node[below] {HTML} ([yshift=-0.2cm]server.west);
                \draw [->, dashed] (server.south) -- (disk.north);
            \end{tikzpicture}
        \end{column}
    \end{columns}
\end{frame}

\begin{frame}{Il Web Dinamico}
    \begin{itemize}
        \item La pagina viene generata \textbf{"al volo"} dal server.
        \item Coinvolge codice (Python, PHP, Node) e Database.
    \end{itemize}
    \centering
    \begin{tikzpicture}[node distance=2cm, scale=0.8, transform shape]
        \node (client) [draw, rounded corners, fill=primary!10, minimum width=2cm] {Client};
        \node (server) [draw, rounded corners, fill=primary!20, right=of client, minimum width=2cm] {App Server};
        \node (db) [draw, shape=cylinder, aspect=0.5, rotate=90, fill=primary!5, right=of server] {\rotatebox{-90}{DB}};
        \draw [->, >=stealth, primary] ([yshift=0.2cm]client.east) -- (server.west);
        \draw [<->] (server.east) -- (db.north);
        \draw [<-, >=stealth, primary] ([yshift=-0.2cm]client.east) -- (server.west);
    \end{tikzpicture}
    \vspace{0.5cm}
    \begin{itemize}
        \item \textbf{Vantaggi:} Personalizzazione, autenticazione, dati in tempo reale.
        \item \textbf{Svantaggi:} Complessità, latenza, sicurezza (SQL Injection).
    \end{itemize}
\end{frame}

% --- CDN ---
\begin{frame}{Cloudflare Pages e CDN}
    \textbf{Perché una CDN?}
    \begin{itemize}
        \item Distribuisce il sito su centinaia di server globali (\textbf{Edge Nodes}).
        \item L'utente scarica dal server fisicamente più vicino.
    \end{itemize}
    \centering
    \begin{tikzpicture}[scale=0.7, transform shape]
        \node (origin) [draw, fill=primary!20, minimum width=2cm] {Origine};
        \node (e1) [circle, draw, fill=primary!5, below left=0.8cm and 1.5cm of origin] {UK};
        \node (e2) [circle, draw, fill=primary!5, below=1cm of origin] {IT};
        \node (e3) [circle, draw, fill=primary!5, below right=0.8cm and 1.5cm of origin] {JP};
        \draw [->, dashed] (origin) -- (e1); \draw [->, dashed] (origin) -- (e2); \draw [->, dashed] (origin) -- (e3);
        \node (u) [draw, rounded corners, fill=green!10, below=0.5cm of e3] {Utente};
        \draw [<->, primary, line width=1pt] (e3) -- (u);
    \end{tikzpicture}
\end{frame}

% --- Rendering ---
\section{Rendering Paradigms}
\begin{frame}{SSR vs CSR}
    \centering
    \begin{tikzpicture}[node distance=3cm, scale=0.8, transform shape]
        \node (s_server) [draw, rounded corners, fill=primary!20, minimum width=2.5cm] {\textbf{SSR}};
        \node (s_client) [draw, rounded corners, fill=primary!10, right=of s_server, minimum width=2.5cm] {HTML Completo};
        \draw [->, >=stealth, primary] (s_server) -- (s_client);
        
        \node (c_server) [draw, rounded corners, fill=primary!20, below=1cm of s_server, minimum width=2.5cm] {\textbf{CSR}};
        \node (c_client) [draw, rounded corners, fill=primary!10, right=of c_server, minimum width=2.5cm] {HTML Vuoto + JS};
        \draw [->, >=stealth, primary] (c_server) -- (c_client);
    \end{tikzpicture}
    \vspace{0.5cm}
    \begin{itemize}
        \item \textbf{SSR (Server-Side)}: Meglio per SEO, caricamento iniziale veloce.
        \item \textbf{CSR (Client-Side)}: Esperienza fluida (SPA), meno carico sul server dopo il primo avvio.
    \end{itemize}
\end{frame}

\begin{frame}{Focus SSR: Il Mondo Django}
    \textbf{Perché Django (Python)?}
    \begin{itemize}
        \item Framework "Batteries Included": gestisce DB, Admin, e Sicurezza.
        \item \textbf{Il Flusso SSR:}
    \end{itemize}
    \centering
    \begin{tikzpicture}[node distance=1.5cm, scale=0.7, transform shape]
        \node (user) [draw, fill=gray!10] {Richiesta URL};
        \node (view) [draw, fill=primary!20, right=of user] {Vista (Python)};
        \node (db) [draw, shape=cylinder, aspect=0.5, rotate=90, fill=primary!5, above=of view] {\rotatebox{-90}{DB}};
        \node (temp) [draw, fill=primary!10, right=of view] {Template (HTML)};
        \node (res) [draw, fill=green!10, below=of temp] {HTML Finale};
        
        \draw [->] (user) -- (view);
        \draw [<->] (view) -- (db);
        \draw [->] (view) -- (temp);
        \draw [->] (temp) -- (res);
    \end{tikzpicture}
    \vspace{0.3cm}
    \begin{itemize}
        \item Il programmatore scrive logica in Python, il server "stampa" l'HTML.
    \end{itemize}
\end{frame}

\begin{frame}[fragile]{Django: Logica nel Template}
    \begin{lstlisting}[language=Python]
<!-- template.html -->
{% if task_list %}
  <ul>
    {% for task in task_list %}
      <li class="{% if task.done %}completed{% endif %}">
        {{ task.name|upper }}
      </li>
    {% endfor %}
  </ul>
{% else %}
  <p>Nulla da fare oggi!</p>
{% endif %}
    \end{lstlisting}
    \begin{itemize}
        \item \textbf{Filtri}: \texttt{|upper} trasforma il testo.
        \item \textbf{Tag}: \texttt{\{\% ... \%\}} per la logica (cicli e condizioni).
        \item \textbf{Variabili}: \texttt{\{\{ ... \}\}} per i dati reali.
    \end{itemize}
\end{frame}

\begin{frame}{Focus SSG: La Velocità di Hugo}
    \textbf{Hugo} è un "SSR eseguito in anticipo":
    \begin{itemize}
        \item \textbf{Zero Database}: I contenuti sono file di testo (Markdown).
        \item \textbf{Build Step}: Hugo trasforma tutto in HTML fisso.
    \end{itemize}
    \vspace{0.3cm}
    \centering
    \begin{tikzpicture}[node distance=1.8cm, scale=0.8, transform shape]
        \node (md) [draw, fill=gray!10] {Post.md};
        \node (hugo) [draw, fill=orange!20, right=of md] {\textbf{HUGO (Build)}};
        \node (html) [draw, fill=green!10, right=of hugo] {index.html};
        
        \draw [->, line width=1.5pt] (md) -- (hugo);
        \draw [->, line width=1.5pt] (hugo) -- (html);
    \end{tikzpicture}
    \vspace{0.3cm}
    \begin{itemize}
        \item \textbf{Performance}: Genera 10.000 pagine in meno di 1 secondo.
        \item \textbf{Deployment}: Perfetto per GitHub Pages o Cloudflare Pages.
    \end{itemize}
\end{frame}

\begin{frame}[fragile]{Hugo: Markdown + Metadati}
    \begin{lstlisting}
---
title: "Il mio nuovo post"
date: 2026-05-14
tags: ["tech", "web"]
---
# Benvenuti

Oggi parliamo di Hugo. Ecco una lista:
- Veloce
- Sicuro
- Gratuito
    \end{lstlisting}
    \begin{itemize}
        \item La parte tra `---` è il \textbf{Frontmatter} (metadati YAML).
        \item Il resto è \textbf{Markdown} (formattazione semplice).
    \end{itemize}
\end{frame}

% --- Lab ---
\section{Laboratorio}
\begin{frame}{Laboratorio: Todo App}
    \textbf{Le fasi del progetto:}
    \begin{enumerate}
        \item \textbf{Struttura (HTML)}: Form e lista.
        \item \textbf{jQuery}: Manipolazione del DOM semplificata.
        \item \textbf{CRUD}: Create, Read, Update, Delete.
        \item \textbf{Persistenza}: Uso del \texttt{localStorage}.
    \end{enumerate}
    \vspace{0.3cm}
    \textbf{Perché jQuery?}
    \begin{itemize}
        \item Sintassi concisa: \texttt{\$(selettore).azione()}.
        \item Astrazione cross-browser.
        \item Facilità di apprendimento per i concetti base.
    \end{itemize}
\end{frame}

\begin{frame}[fragile]{Persistenza con localStorage}
    \begin{lstlisting}[language=Java]
// Salvare i dati (convertendo in stringa)
localStorage.setItem('tasks', JSON.stringify(tasks));

// Recuperare i dati (convertendo in oggetto)
const data = JSON.parse(localStorage.getItem('tasks'));
    \end{lstlisting}
    \vspace{0.3cm}
    \textbf{Nota:} Il localStorage sopravvive al ricaricamento della pagina e alla chiusura del browser!
\end{frame}

\begin{frame}
    \centering
    \Huge \textcolor{primary}{Iniziamo il Laboratorio!}
    \vspace{1cm}
    \normalsize
    Aprite il file \texttt{index.html} nella cartella \texttt{src/lezione2/3.todo-vanilla}
\end{frame}

\end{document}
